
\section{Conclusion}

\subsection{Summary}

We presented \textbf{A3} (Activation-Aware Adaptive factorization), a low-rank 
compression method that combines Fisher-diagonal weighting, streaming 
calibration, strict parameter rank ceilings ($r_{\mathrm{ceiling}} \le 256$), and closed-form LayerNorm recalibration. A3 
guarantees positive parameter savings across all architectures (+1.25M parameters saved on DistilBERT, 
+0.59M on RoBERTa, +0.53M on GPT-2) while preserving task performance across SST-2, MNLI, and AG News.

\subsection{Results Recap}

\begin{table}[H]
\centering
\small
\begin{tabular}{lllll}
\toprule
\textbf{Model} & \textbf{Validation Acc} & \textbf{Acc $\Delta$} & \textbf{Params Saved} & \textbf{Peak VRAM} \\
\midrule
DistilBERT & 90.02\% & -0.69\% & +1,245,696 & 334 MB \\
RoBERTa    & 83.26\% & +4.70\% & +589,824   & 531 MB \\
GPT-2      & 63.19\% & -16.17\% & +525,312  & 572 MB \\
\bottomrule
\end{tabular}
\caption{Consolidated experimental results across the three target architectures on SST-2.}
\label{tab:conclusion_recap}
\end{table}
% <!-- source: results/a3_tests/fixed_v4/distilbert_sst2_a3_full_metrics.json line 4 -->
% <!-- source: results/unified_a3/roberta/roberta-base_a3_full_metrics.json line 4 -->
% <!-- source: results/unified_a3/gpt2/gpt2_a3_full_metrics.json line 4 -->

\subsection{Key Takeaways}

\begin{enumerate}
\item \textbf{Guaranteed Parameter Reduction:} Enforcing a strict rank ceiling below the break-even dimension eliminates parameter expansion, ensuring every factorized layer saves memory.
\item \textbf{Fisher Weighting is Essential:} Guiding the SVD input space with empirical diagonal sensitivity retains critical feature directions, contributing a +3.6\% to +10.1\% accuracy gain over standard SVD.
\item \textbf{LayerNorm Recalibration:} Closed-form post-factorization affine scale correction stabilizes activation variance, closing over 60\% of initial representation shrinkage before fine-tuning.
\item \textbf{Safe Fallback Mechanism:} Preserving original dense weights on high-entropy layers (such as deep value and output projections) safeguards against catastrophic model degradation.
\end{enumerate}
\subsection{Future Work}

\subsubsection{Short-Term}

\begin{itemize}
\item \textbf{Fused Triton kernels} for factorized forward pass
\item \textbf{Multi-task calibration} for broader generalization
\item \textbf{Combined with quantization} for 8$\times$ compression
\item \textbf{Rank schedule learning} via Bayesian optimization

\end{itemize}
\subsubsection{Long-Term}

\begin{itemize}
\item \textbf{Task-agnostic calibration:} Use LLM-generated text
\item \textbf{Recursive factorization:} Apply A3 to factorized layers
\item \textbf{Theoretical analysis:} Bound error propagation across layers

\end{itemize}
\subsection{Closing Remarks}

A3 is a step toward practical, honest model compression. We release all code 
and checkpoints to enable replication and incremental improvement.

\subsection{Acknowledgments}

We thank the anonymous reviewers and the compute cluster administrators for their feedback and support. This research was supported by the Advanced Machine Learning \& Efficiency Laboratory.

\subsection{Author Contributions}

\begin{itemize}
\item \textbf{Taha Majs} designed A3, implemented the code, ran experiments, wrote
  the paper.
\item \textbf{Collaborative Research Team} provided methodological feedback and edited the manuscript.

\end{itemize}
\subsection{Conflict of Interest}

The authors declare no competing interests.